\documentclass{article}
\usepackage[utf8]{inputenc}
\usepackage[T1]{fontenc}
\usepackage{geometry}
\usepackage{graphicx}
\usepackage{float}
\usepackage{hyperref}
\usepackage[english,indonesian]{babel}

\title{Optimasi Pemrosesan Template LaTeX Menggunakan Pendekatan Berbasis AI}
\author{Academic Support Agent}
\date{\today}

\begin{document}
\maketitle

\begin{abstract}

\end{abstract}

\chapter{PENDAHULUAN}


\chapter{TINJAUAN PUSTAKA}


% Original content (commented out):
% % Modele de doc LaTeX pour bac blanc  avec page de garde style bac et préambule de l'APMEP adapté pour Overleaf et mise en page A4.
%http://www.apmep.fr/-Annales-Bac-Brevet-BTS-
%http://www.apmep.fr/-Terminale-S-240-sujets-depuis-1999-
% 
% Author: Vincent Pantaloni (Lycée Jean Zay, Orleans)
%  http://prof.pantaloni.free.fr/
%  Twitter @panlepan
\documentclass[a4paper,10pt]{article}
\usepackage[T1]{fontenc}
%\usepackage[utf8]{inputenc} %commenté pour Overleaf
%\usepackage{fontspec}% Warning Overleaf à laisser commenté.
\usepackage{fourier}
\usepackage[scaled=0.875]{helvet} 
\renewcommand{\ttdefault}{lmtt} 
\usepackage{amsmath,amssymb,makeidx}
\usepackage[normalem]{ulem}
\usepackage{fancybox,graphicx}
\usepackage{tabularx}
\usepackage{dcolumn}
\usepackage{textcomp}
\usepackage{diagbox}
\usepackage{tabularx}
\usepackage{lscape}
\newcommand{\euro}{\eurologo{}}
\usepackage{pstricks,pst-plot,pst-text,pst-tree,pstricks-add}
%Modif mise en page de APMEP:
%\setlength\paperheight{297mm}
%\setlength\paperwidth{210mm}
%\setlength{\textheight}{25cm}
%Remplacé par geometry:
\usepackage[marginratio=1:1,textwidth=17cm,textheight=23cm,headheight=2cm,columnsep=0.7cm]{geometry}
\newcommand{\R}{\mathbb{R}}
\newcommand{\N}{\mathbb{N}}
\newcommand{\D}{\mathbb{D}}
\newcommand{\Z}{\mathbb{Z}}
\newcommand{\Q}{\mathbb{Q}}
\newcommand{\C}{\mathbb{C}}
\renewcommand{\theenumi}{\textbf{\arabic{enumi}}}
\renewcommand{\labelenumi}{\textbf{\theenumi.}}
\renewcommand{\theenumii}{\textbf{\alph{enumii}}}
\renewcommand{\labelenumii}{\textbf{\theenumii.}}
\newcommand{\vect}[1]{\mathchoice%
{\overrightarrow{\displaystyle\mathstrut#1\,\,}}%
{\overrightarrow{\textstyle\mathstrut#1\,\,}}%
{\overrightarrow{\scriptstyle\mathstrut#1\,\,}}%
{\overrightarrow{\scriptscriptstyle\mathstrut#1\,\,}}}
\def\Oij{$\left(\text{O},~\vect{\imath},~\vect{\jmath}\right)$}
\def\Oijk{$\left(\text{O},~\vect{\imath},~\vect{\jmath},~\vect{k}\right)$}
\def\Ouv{$\left(\text{O},~\vect{u},~\vect{v}\right)$}
\setlength{\voffset}{-1,5cm}
\usepackage{fancyhdr}
\usepackage{hyperref}

%\hypersetup{%pdfauthor = {APMEP},pdfsubject = {Baccalauréat S},pdftitle = {Métropole -- La Réunion 20 juin 2016},allbordercolors = white,pdfstartview=FitH}   

\usepackage[frenchb]{babel}
\usepackage[np]{numprint}
\setlength\parindent{0mm}%Pas d'indentation  après saut de ligne.
%% Ajouts pour page de garde style BAC
\usepackage[firstpage]{draftwatermark}%

\begin{document}
\renewcommand{\labelitemi}{\textbullet} %pour éviter les tirets dans les "itemize" qui apportent confusion avec le signe moins.
%% Début page de garde style BAC
\SetWatermarkText{{\bf O B L I G A T O I R E}}
\SetWatermarkFontSize{1.8cm}
\SetWatermarkLightness{0.5}
\pagestyle{fancy}
\lhead{{\sl Lycée Jean Zay}}%\textbf{Nom:}}
\chead{}%\textbf{Contr\^ole commun \no 3}}
\rhead{\textbf{Baccalauréat blanc}$\ $ page \thepage / 5 }
\lfoot{} \cfoot{}\rfoot{}
\rule{0pt}{1.5cm}
\begin{center}
\fbox{{\Huge BACCALAUR\'EAT G\'EN\'ERAL}}\\
\vspace{3cm}
Session 2017\\
\vspace{2cm}

{\LARGE {\bf MATH\'EMATIQUES}}\\
\rule{7cm}{1pt}\\
\vspace{1.5cm}
{\Large Série: {\bf S}}\\
\vspace{1.5cm}
DUR\'EE DE L'\'EPREUVE: {\bf 4 heures} -- COEFFICIENT: {\bf 7}\\
\vspace{2cm}
{\bf {\it Ce sujet comporte cinq pages numérotées de 1/5 à 5/5}}\\
{\bf {\it La dernière page comporte une annexe à remettre complétée de
    votre nom et classe.}}\\
\vspace{2cm}
{\it L'utilisation d'une calculatrice est autorisée}\\
\vspace{1.5cm}
\framebox[1.1\width]{
\parbox{0.78\textwidth}{La qualité de la rédaction et la précision des raisonnements  entreront pour une part importante dans l'appréciation des copies.}}
\end{center}
\vfill
%\vspace{3cm}
\hfill {\bf Tournez la page S.V.P.}
%% FIN page de garde style BAC
\newpage
%
%
%

%
%%###########################################
%%  DEBUT --- COLLAGE ICI DES EXOS APMEP
%%###########################################
%
%%%%%%%%%%%%%%%%  EXO 1
%Métropole -- La Réunion 20 juin 2016
\textbf{\textsc{Exercice 1 \hfill 6 points}}

\textbf{Commun à tous les candidats} 

\medskip

Les trois parties sont indépendantes. Les résultats des probabilités seront arrondis à $10^{-3} $ près.

\medskip

\textbf{Partie 1}

\medskip

On estime qu'en 2013 la population mondiale est composée de 4,6 milliards de personnes âgées de 20 à 79 ans et que 46,1\,\% des personnes âgées de 20 à 79 ans vivent en zone rurale et 53,9\,\% en zone urbaine.

En 2013, d'après la fédération internationale du diabète, 9,9\,\% de la population mondiale âgée de 20 à 79 ans vivant en zone urbaine est atteinte de diabète et 6,4\,\% de la population mondiale âgée de 20 à 79 ans vivant en zone rurale est atteinte de diabète.

On interroge au hasard une personne âgée de 20 à 79 ans. On note :

$R$ l'évènement : \og la personne choisie habite en zone rurale \fg,

$D$ l'évènement: \og la personne choisie est atteinte de diabète \fg.

\medskip

\begin{enumerate}
\item Traduire cette situation à l'aide d'un arbre de probabilité.
\item 
	\begin{enumerate}
		\item Calculer la probabilité que la personne interrogée soit diabétique.
		\item La personne choisie est diabétique. Quelle est la probabilité qu'elle habite en zone rurale ?
 	\end{enumerate}
\end{enumerate}

\bigskip

\textbf{Partie 2}

\medskip

Une personne est dite en hypoglycémie si sa glycémie à jeun est inférieure à 60~mg.dL$^{-1}$ et elle est en hyperglycémie si sa glycémie à jeun est supérieure à 110~mg. dL$^{-1}$. La glycémie à jeun est considérée comme \og normale \fg{} si elle est comprise entre 70~mg. dL$^{-1}$ et 110 mg.dL$^{-1}$. Les personnes ayant un taux de glycémie compris entre 60 et 70~mg.rdL$^{-1}$ ne font pas l'objet d'un suivi particulier.

On choisit au hasard un adulte dans cette population. Une étude a permis d'établir que la probabilité qu'il soit en hyperglycémie est 0,052 à $10^{-3}$ près. Dans la suite on admettra que cette probabilité est égale à $0,052$.

On modélise la glycémie à jeun, exprimée en mg.dL$^{-1}$, d'un adulte d'une population donnée, par une variable aléatoire $X$ qui suit une loi normale d'espérance $\mu$ et d'écart-type $\sigma$.

On donne ci-dessous la représentation graphique de la densité de probabilité de la variable aléatoire $X$.


\begin{center}
\psset{xunit=0.1cm, yunit=40cm, arrowsize=2pt 3, algebraic=true}
\def\xmin {45}
\def\xmax {135}
\def\ymin {-0.02}
\def\ymax {0.035}
\begin{pspicture*}(\xmin,\ymin)(\xmax,\ymax)
\def\m{90} \def\s{12.3}
\def\f{1/(\s*sqrt(2*PI))*EXP((-((x-\m)/\s)^2)/2)}

\psplot[plotpoints=1000]{\xmin}{\xmax}{\f}

\uput[d](70,0){70}
\uput[d](90,0){$\mu = 90$}
\uput[d](110,0){110}
\uput[d](125,0){mg.dL$^{-1}$}

\psline[linestyle=dashed](90,0)(90,0.0324)
\psline[linestyle=dashed](70,0)(70,0.00865)
\psline[linestyle=dashed](110,0)(110,0.00865)
\psline(\xmin,0)(\xmax,0)
\end{pspicture*}
\end{center}


\begin{enumerate}
\item Quelle est la probabilité que la personne choisie ait une glycémie à jeun \og normale\fg{} ?
\item Déterminer la valeur de $\sigma$ arrondie au dixième.
\item Dans cette question, on prend $\sigma = 12$. Calculer la probabilité que la personne choisie soit en hypoglycémie.
\end{enumerate}
 
\bigskip
 
\textbf{Partie 3}
 
 \medskip
 
Afin d'estimer la proportion, pour l'année 2013, de personnes diagnostiquées diabétiques dans la population française âgée de 20 à 79 ans, on interroge au hasard \np{10000}~personnes.
 
Dans l'échantillon étudié, 716 personnes ont été diagnostiquées diabétiques.
 
\medskip
 
\begin{enumerate}
\item À l'aide d'un intervalle de confiance au niveau de confiance 95\,\%, estimer la proportion de personnes diagnostiquées diabétiques dans la population française âgée de 20 à 79 ans.
\item Quel doit être le nombre minimal de personnes à interroger si l'on veut obtenir un intervalle de confiance d'amplitude inférieure ou égale à 0,01 ?
\end{enumerate}

%\vspace{0,5cm}
\newpage
%%%%%%%%%%%%%%%%  EXO 2
%Métropole -- La Réunion 20 juin 2016
\textbf{\textsc{Exercice 2 \hfill 4 points}}

\textbf{Commun à tous les candidats} 

\medskip

On considère les nombres complexes $z_n$ définis pour tout entier $n \geqslant 0$ par la donnée de $z_0$, où $z_0$ est différent de 0 et de 1, et la relation de récurrence:

\[z_{n+1} = 1- \dfrac{1}{z_n}.\]

\begin{enumerate}
\item 
	\begin{enumerate}
		\item Dans cette question, on suppose que $z_0 = 2$. Déterminer les nombres $z_1$, $z_2$, $z_3$,  $z_4$,  $z_5$ et $z_6$
		\item Dans cette question, on suppose que $z_0 = \text{i}$. Déterminer la forme algébrique des nombres complexes $z_1$, $z_2$, $z_3$,  $z_4$,  $z_5$ et $z_6$.
\item Dans cette question on revient au cas général où $z_0$ est un complexe donné. Que peut-on conjecturer pour les valeurs prises par $z_{3n}$ selon les valeurs de l'entier naturel $n$ ?

Prouver cette conjecture.
	\end{enumerate}
\item Déterminer $z_{\np{2016}}$ dans le cas où $z_0 = 1 + \text{i}$.
\item Existe-t-il des valeurs de $z_0$ tel que $z_0 = z_1$ ? Que peut-on dire de la suite $\left(z_n\right)$ dans ce cas ?
\end{enumerate}
%%%%%%%%%%%%%%%%  EXO 3
\vspace{0,5cm}
\textbf{\textsc{Exercice 3 \hfill 5 points}}

\textbf{Candidats n'ayant pas suivi l'enseignement de spécialité}
\end{document}


\end{document}